\subsection{Authenticated Encryption}

So far we have considered security and 
integrity seperately. Encryption thus 
far has not guaranteed integrity, and 
message authentication codes do not 
guarantee confidentiality. In fact, 
in many a scenario, lack of integrity 
implies lack of confidentiality. 

This leads to the need for 
\emph{authenticated encryption}, 
which incorporates a checksum. 
A simple way to do this is tag 
the message with a checksum and 
then encrypt both the message 
and tag (auth-then-enc). You 
may also send the MAC for the 
plaintext with the encrypted message 
(auth-and-enc). It turns out the 
only secure method is to encrypt 
and then compute the MAC for the 
encrypted message. This method, 
enc-then-auth, is secure. 
There are still several vulnerabilities:
\begin{itemize}
    \item Re-ordering attack: swap order of messages
    \item Replay attack: replay a valid ciphertext previously sent
    \item Message dropping attack: drop some of the messages sent
    \item Reflection attack: take a 
    ciphertext from Alice to Bob and 
    send it back to Alice
\end{itemize}
To defeat these vulnerabilities, keep two counters 
\texttt{ctr}$_{A,B}$, \texttt{ctr}$_{B,C}$ and 
define a directionality bit $b$. A sends 
$c = \text{Enc}_k(b_{A,B}|\text{\texttt{ctr}}_{A,B}|m)$.